\usepackage[mathscr]{eucal}
\usepackage[dvipsnames]{xcolor}

\def\violet {\color[rgb]{1,.4,.7}}
\def\red{\color[rgb]{1,0,0}}
\def\verde{\color[rgb]{.4,.4,0}}
\def\blu{\color[rgb]{0,0,1}}
\def\gris{\color[rgb]{.4,0.4,0.4}}
\def\bla{\color[rgb]{0,0,0}}
\def\ora{\color[rgb]{1,.6,0}}
\def\gre{\color[rgb]{.4,.4,.4}}
\def\colr{\color[rgb]{.4,0,0}}
\def\colb{\color[rgb]{0,0,.4}}
\definecolor{red}{rgb}{1,0,0}
\definecolor{purple}{rgb}{0.63, 0.36, 0.94}
%\definecolor{darkgreen}{rgb}{0,100,0}
\definecolor{Kelly}{rgb}{.1,.65,.1}

\def\Xint#1{\mathchoice
{\XXint\displaystyle\textstyle{#1}}%
{\XXint\textstyle\scriptstyle{#1}}%
{\XXint\scriptstyle\scriptscriptstyle{#1}}%
{\XXint\scriptscriptstyle\scriptscriptstyle{#1}}%
\!\int}
\def\XXint#1#2#3{{\setbox0=\hbox{$#1{#2#3}{\int}$ }
\vcenter{\hbox{$#2#3$ }}\kern-.6\wd0}}
\def\ddashint{\Xint=}
\def\dashint{\Xint-}

\DeclareFontFamily{U}{mathx}{}
\DeclareFontShape{U}{mathx}{m}{n}{<-> mathx10}{}
\DeclareSymbolFont{mathx}{U}{mathx}{m}{n}
\DeclareMathAccent{\widehat}{0}{mathx}{"70}
\DeclareMathAccent{\widecheck}{0}{mathx}{"71}

\def\opE{\mathfrak E} % originally \mathcal
\def\opD{\mathfrak D}
\def\opR{\mathfrak R}
\def\SU{{\rm SU}}
\def\SO{{\rm SO}}
\def\rS{{\mathrm S}}
\def\rs{{\mathrm s}}
\newcommand\Leg{ {\mathcal P}} % [closed] regular parametrized Legendrian curves (add subscript for periodic)
\newcommand\Legq{\mathcal Q}
\newcommand\Legu{ {\mathcal P}^1} % [closed] unit-speed curves Legendrian curves (add subscript for length)
\newcommand\Leguq{{\mathcal Q}^1}

\def\Dx{\delta}		% total x-derivative
\def\Fop{\mathtt F} %Kind of ugly
\def\Euler{\mathtt E}
\def\Phit{\widetilde\Phi}

\newcommand\uM{\mathsf M} %RHS of mKdV hierarchy
\newcommand\uL{\mathsf L} % antiderivative of M
\newcommand\uN{\mathsf N} % antiderivative of u M
\newcommand\uP{\mathsf P} %generic differential polynomial
\newcommand\uQ{\mathsf Q} % ditto
\makeatletter % mKdV dependent variable and its derivatives
\newcommand{\uu}[1][\@nil]{%
 \def\tmp{#1}%
 \ifx\tmp\@nnil
 {u}
 \else
 {u}_{#1}
 \fi}
\makeatother

\newcommand\kk{k} %Legendrian curvature -- could use \kappa instead
\newcommand\speed{\beta} %speed of a Legendrian curve
\def\ri{\mathrm{i}}
\def\vecV{\mathbf V}
\def\vecW{\mathbf W}
\newcommand\VV{\mathbf V}
\newcommand\VW{\mathbf W}
\newcommand\VR{\mathbf R}
\newcommand\VH{\mathbf H}
\newcommand\VZ{\mathbf Z}
\newcommand\qV{\mathcal V}%vector fields on the quotient
\newcommand\qW{\mathcal W}
\newcommand\calC{\mathcal C} %sets in the cnoidal moduli space
\newcommand\funH{{\mathscr H}} %functionals on the space of Legendrian curves
\newcommand\funA{{\mathscr A}} %total arclength

\newcommand{\dx}{\,{\rm d}x}
\newcommand{\ds}{\,{\rm d}s}
\newcommand{\sdown}{\check{s}} %arclength parameter along Clifford projection
\def\gam{\gamma} % usually gamma is too low, but don't use this one in subscripts
\def\di{\partial}
\def\C{\mathbb C}
\def\R{\mathbb R}
\def\Z{\mathbb Z}
\newcommand\rU{\mathrm U}
\newcommand\bw{\mathbf w}
\newcommand\defem{\sl} %emphasis used in definitions -- could be boldface instead
\newcommand\fsu{\mathfrak{su}}
\newcommand\sj{\mathsf j}
\newcommand\Ad{\operatorname{Ad}}
\newcommand\cl{\mathbf{cl}}
\newcommand\ve{\mathbf e} %not sure
\newcommand\be{{\bf e}}
\newcommand\bp{{\bf p}}
\newcommand\bph{{\bp_H}} %Heisenberg projection
\newcommand\gamt{\widetilde{\gam}}
\newcommand\gamf{\widehat{\gam}} % for a family of curves, in variational context
\newcommand\Gamf{\widehat{\Gamma}} % family of Frenet frames
\newcommand\kh{\widehat{\kk}} % for a family of curves
\newcommand\bert{\big\vert}
\newcommand\rd{\mathrm{d}}
\newcommand\rA{\mathrm{A}}
\newcommand\rC{\mathrm{C}}
\newcommand\rK{\mathrm{K}}
\newcommand\tJ{\mathtt{J}}
\newcommand\tU{\mathtt{U}}
\newcommand\tV{\mathtt{V}}
\renewcommand\Re{\operatorname{Re}}
\renewcommand\Im{\operatorname{Im}}
\newcommand\cn{\operatorname{cn}}
\newcommand\dn{\operatorname{dn}}
\newcommand\aside[1]{{\tt\color{blue}{#1}}}

\newcommand\annac[1]{{\tt\color{purple}#1}}% Anna's comments
\newcommand\annaTag[1]{\color{brown}#1} %edits made by Anna with which Tom agrees
\newcommand\tomedit[1]{{\color{PineGreen}#1}}
\newcommand\tomeditopen {\color{PineGreen}} % for longer passages; must use inside a curly brace
\newcommand\tomcomment[1]{{\tt \color{Kelly}#1}}
\newcommand\emiliocomment[1]{{\tt \color[rgb]{1,.5,0} #1}}

\numberwithin{equation}{section}

\newtheorem{Theorem}{Theorem}[section]
\newtheorem*{Theorem*}{Theorem}
\newtheorem{Corollary}[Theorem]{Corollary}
\newtheorem{Lemma}[Theorem]{Lemma}
\newtheorem{Proposition}[Theorem]{Proposition}
 { \theoremstyle{definition}
\newtheorem{Definition}[Theorem]{Definition}
\newtheorem{Note}[Theorem]{Note}
\newtheorem{Example}[Theorem]{Example}
\newtheorem{Remark}[Theorem]{Remark} }
\def\gami{\gam} % source line 925 verbatim (local notation macro used by the timeevolution span)
